\documentclass{article}
\usepackage{amsmath, amsthm, amsfonts, amssymb}
\theoremstyle{plain}
\newtheorem{thm}{Theorem}
\numberwithin{thm}{section}

\theoremstyle{plain}
\newtheorem{prop}{Proposition}
\numberwithin{prop}{section}

\theoremstyle{definition}
\newtheorem{defn}{Definition}
\numberwithin{defn}{section}

\theoremstyle{remark}
\newtheorem*{rem}{Remark}

\newtheorem*{cor}{Corollary}

\numberwithin{equation}{section}

\newcommand{\op}{\prime}
\begin{document}
\title{Linear Analysis}
\author{Amey Joshi}
\maketitle
\section{The geometry of function spaces}\label{s1}
We start with the definition of a vector space over a field.
\begin{defn}\label{s1d1}
A set $V$ is called a vector space over a field $\mathbb{F}$ if:
\begin{enumerate}
\item $V$ is an abelian group with the operation $+$.
\item Scalar multiplication is defined as a mapping $\mathbb{F} \times V \mapsto
V$ with the properties $\alpha(x + y) = \alpha x + \alpha y$ and $(\alpha +
\beta) x = \alpha x + \beta x$ for all $\alpha, \beta \in \mathbb{F}$ and $x, y
\in V$.
\end{enumerate}
\end{defn}

Let $\Omega = [a, b]$ and $C^k(\Omega)$ be the set of functions $f: \Omega 
\rightarrow \mathbb{R}$ such that first $k$ derivatives of $f$ are continuous
in $\Omega$. We observe that
\begin{prop}\label{s1p1}
$C^1(\Omega)$ is a vector space over $\mathbb{R}$.
\end{prop}
\begin{proof}
The set of all real-valued functions on $\Omega$ is a vector space. To prove
that its subset $C^1{\Omega}$ is also a vector space, it suffices to show that
for all $f, g \in C^1(\Omega), f + g \in C^1(\Omega)$ and for all $\alpha \in
\mathbb{R}, f \in C^1(\Omega), \alpha f \in C^1(\Omega)$.

$f \in C^1(\Omega)$ implies that $Df$ exist and is continuous in $\Omega$. 
Therefore, for a given $\epsilon > 0$, we can find $\delta_1 > 0$ such that
$|x - x_0| < \delta_1 \Rightarrow |Df(x) - Df(x_0)| < \epsilon/2$. Similarly,
we can find $\delta_2 > 0$ such that $|x - x_0| < \delta_2 \Rightarrow |Dg(x)
- Dg(x_0) < \epsilon/2$. Choose $\delta = \min(\delta_1, \delta_2)$ for all 
$|x - x_0| < \delta, |D(f + g)(x) - D(f + g)(x_0)| = |Df(x) - Df(x_0) + Dg(x) - 
Dg(x_0)| \le |Df(x) - Df(x_0)| + |Dg(x) - Dg(x_0)| < \epsilon$. Thus $f + g \in
C^1(\Omega)$. 

Let $\alpha \ne 0$. Find $\delta > 0$ such that $|x - x_0| < \delta \Rightarrow
|Df(x) - Df(x_0)| < \epsilon/|\alpha|$. That is $|\alpha||Df(x) - Df(x_0)| < 
\epsilon \Rightarrow |D(\alpha f)(x) - D(\alpha f)(x_0)| < \epsilon$. Thus,
$\alpha f \in C^1(\Omega)$. 
\end{proof}

\begin{defn}[Linear independence]\label{s1d2}
A subset $\{v_1, \ldots, v_n\}$ of $V$ is said to be linear independent if 
$\alpha_1 v_1 + \cdots + \alpha_n v_n = 0 \Leftrightarrow \alpha_1 = 0, \ldots,
\alpha_n = 0$.
\end{defn}

\begin{defn}[Linear combination]\label{s1d3}
A linear combination of a subset $\{v_1, \ldots, v_n\}$ of $V$ is the set $\{
\alpha_1 v_1 + \cdots + \alpha_n v_n\}$ for any $\alpha_1, \ldots, \alpha_n \in
\mathbb{F}$.
\end{defn}

\begin{defn}[Span]\label{s1d4}
The span of a subset $\{v_1, \ldots, v_n\}$ of $V$ is the set of all its linear
combinations.
\end{defn}

\begin{defn}[Basis]\label{s1d5}
A subset $\{v_1, \ldots, v_n\}$ of $V$ is called a basis of $V$ is its span is
$V$. That is, any member of $V$ can be written as a linear combination of $v_1,
\ldots, v_n$.
\end{defn}

\begin{defn}[Dimension]\label{s1d6}
The dimension of a vector space is the number of elements in its basis.
\end{defn}

\begin{prop}[Steinitz exchange lemma]\label{s1p2}
Let $U$ and $W$ be finite subsets of a vector space $V$. If $U$ is a set of 
linearly independent vectors and $W$ spans $V$ then:
\begin{enumerate}
\item $|U| \le |W|$ and
\item There is a set $W^\op \subseteq W$ such that $|W^\op| = |W| - |U|$ such
that $U \cup W^\op$ spans $V$.
\end{enumerate}
\end{prop}
\begin{proof}
Let $U = \{u_1, \ldots, u_m\}$ and $W = \{w_1, \ldots, w_n\}$. We will show 
that $m \le n$ and that after rearranging $w_i$'s if needed, the set $\{u_1,
\ldots, u_m, w_{m+1}, \ldots w_n\}$ spans $V$. The proof is by induction on $m$.

If $m = 0$, $U$ is empty and $W$ spans $V$. Furthermore, $W^\op = W$. This
proves the base case. Let the claim be true for $m = k-1$. Then $\{u_1, \ldots,
u_{k-1}, w_k, \ldots, w_n\}$ spans $V$. Therefore,
\[
u_k = \sum_{i=1}^{k-1}\alpha_iu_i + \sum_{i=k}^n\alpha_iw_i.
\]
At least one of $\alpha_k, \ldots, \alpha_n$ must be non-zero. Otherwise, $u_k$
will be a linear combination of $u_1, \ldots, u_{k-1}$ contradicting our
hypothesis that $U$ is a set of linearly independent vectors. Without loss of
generality, assume that $\alpha_k \ne 0$. Then we can write
\[
w_k = -\sum_{i=1}^{k-1}\frac{\alpha_i}{\alpha_k} - \frac{1}{\alpha_k}u_k - 
\sum_{i=k+1}^n\alpha_iw_i,
\]
that is the set $\{u_1, \ldots, u_k, w_{k+1}, \ldots, w_n\}$ spans $V$. Thus,
by induction, we can show that $\{u_1, \ldots, u_m, w_{m+1}, \ldots, w_n\}$ 
spans $V$. Furthermore, since the induction step involves building a basis of
size $n$ by replacing one of $w_i$'s with a $u_k$, $U$ cannot have more elements
than $W$. That is $m \le n$. Also the set $W^\op = \{w_{m+1}, \ldots, w_n\}$.
\end{proof}
\begin{rem}
Since the finite set $W$ spans the vector space $V$, this theorem is applicable
only to finite-dimensional spaces.
\end{rem}

\begin{prop}\label{s1p3}
Every basis of $V$ has the same number of elements.
\end{prop}
\begin{proof}
Let, if possible, $V$ have two bases $B_1$ and $B_2$ such that $|B_1| = m$ and
$|B_2| = n$. Since $B_1$ is linearly independent ad $B_2$ spans $V$, by Steinitz
exchange lemma (\ref{s1p2}), $n \le m$. Similarly, since $B_2$ is linearly
independent and $B_1$ spans $V$, $m \le n$.
\end{proof}

We will now show that there are vector spaces of infinite dimensions.
\begin{prop}\label{s1p4}
The set $P(\mathbb{R})$ of all polynomials is an infinite-dimensional vector 
space.
\end{prop}
\begin{proof}
Since a sum of two polynomials is a polynomial and a polynomial multiplied by a
real number is also a polynomial, the set $P(\mathbb{R})$ is a vector space. 
Let, if possible, it have $n$ dimensions. Let $B = \{p_1, \ldots, p_n\}$ be its
basis. For polynomials $p, q$, the degree of their sum is at most the maximum of
degrees of $p$ and $q$ and the degree of $\alpha p$ is smme as that of $p$.
Therefore, if $M$ is the maximum degree of polynomials in this basis, we cannot
construct a polynomial of degree more than $M$ by a linear combination of $p_1,
\ldots, p_n$. Thus, $B$ does not span $P(\mathbb{R})$ and cannot be its basis.
\end{proof}
\begin{rem}
The nature of the coefficients is irrelavant for this property. They may be
real, complex or belonging to a finite field.
\end{rem}

\section{Distance, magnitude and convergence}\label{s2}
\begin{defn}[Norm]\label{s2d1}
The norm on a vector space is a function $\lVert\cdot\rVert:V\mapsto\mathbb{R}$
such that
\begin{enumerate}
\item $\lVert v \rVert \ge 0 \;\forall\; v \in V$ and $\lVert v \rVert = 0
\Leftrightarrow v = 0$.
\item $\lVert \alpha v\rVert = |\alpha|\lVert v \rVert$.
\item $\lVert u + v\rVert \le \lVert u \rVert + \lVert v \rVert\;\forall\; u, v
\in V$.
\end{enumerate}
\end{defn}

Consider the space $C^0[0, 1]$ of continuous functions on $[0, 1]$ and define 
the function $\lVert\cdot\rVert_1$ as
\begin{equation}\label{s2e1}
\lVert f \rVert_1 = \int_0^1 |f(x)|dx.
\end{equation}
\begin{prop}\label{s2p1}
The function defined by \eqref{s2e1} is a norm.
\end{prop}
\begin{proof}
Let $f \in C^0[0, 1]$. Since $f$ is continuous in a bounded interval, it is
bounded. That is, $\exists\;m, M \in \mathbb{R}$ such that $m \le f(x) \le M$
for all $x \in [0, 1]$. Choose $M^\ast = \max(M, -m)$ then $|f(x)| \le M^\ast$
for all $x \in [0, 1]$. Therefore, $\lVert f \rVert|_1 \le M\int_0^1dx = M^\ast
\ge 0$. If $f = 0$ then $\lVert f \rVert_1 = 0$. If $\lVert f \rVert_1 = 0$
then $\int_0^1|f|dx = 0$. Since $f$ is continuous, this implies that $|f| = 0$.

Consider the second property,
\[
\lVert\alpha f\rVert_1 = \int_0^1|\alpha f(x)|dx = |\alpha|\int_0^1|f(x)|dx
= |\alpha|\lVert f \rVert_1.
\]

Finally,
\begin{eqnarray*}
\lVert f + g\rVert_1 &=& \int_0^1|f(x) + g(x)|dx \\
    &\le& \int_0^1|f(x)|dx + \int_0^1|g(x)|dx \\
    &\le& \lVert f \rVert_1 + \lVert g \rVert_1
\end{eqnarray*}    
\end{proof}
\begin{rem}
The norm defined by \eqref{s2e1} is called the $L^1$ norm.
\end{rem}

\begin{defn}[A metric]\label{s2d2}
A metric on a vector space $V$ is a mapping $d: V \times V \mapsto \mathbb{R}$
with the following properties:
\begin{enumerate}
\item $d(x, y) \ge 0$. $d(x, y) = 0 \Leftrightarrow x = y$.
\item $d(x, y) = d(y, x)$,
\item $d(x, z) \le d(x, y) + d(y, z)$,
\end{enumerate}
for all $x, y, z \in V$.
\end{defn}

Every norm introduces a metric $d(x, y) = \lVert x - y\rVert$. Observe that
the addition in the definition of a norm is an operation on the members of
a vector space while that in the definition of a metric is the addition of real
numbers. Thus, a metric is a broader concept than a norm. If a metric is
induced by a norm then it has the following properties.
\begin{prop}\label{s2p2}
Let $V$ be a vector space with a norm $\lVert\cdot\rVert$ and let $d$ be a 
metric on $V$. If $d$ is induced by the norm then,
\begin{enumerate}
\item $d(x + u, y + u) = d(x, y)$, that is $d$ is translation invariant and
\item $d(\alpha x, \alpha y) = |\alpha|d(x, y)$, that is, $d$ is homogeneous.
\end{enumerate}
\end{prop}
\begin{proof}
$d(x + u, y + u) = \lVert x + u - (y + u)\rVert=\lVert x - y \rVert = d(x, y)$.
Likewise, $d(\alpha x, \alpha y) = \lVert\alpha x - \alpha y\rVert = \lVert
\alpha(x - y)\rVert = |\alpha|\lVert x - y\rVert = |\alpha|d(x, y)$.
\end{proof}

Define the discrete metric as
\[
d(x, y) = \begin{cases}
0 \text{ if }x = y \\
1 \text{ otherwise.}
\end{cases}
\]
If $x \ne y$ and $y \ne z$, $d(x, y) + d(y, z) = 2 \ge d(x, z) = 1$. Thus, $d$
is indeed a metric. However, it is not induced by a norm because $d(\alpha x,
\alpha y) = 1 = d(x, y)$ whenever $x \ne y$.

\begin{defn}[Complete space]\label{s2d3}
A $V$ be a vector space with a norm $\lVert\cdot\rVert$ and $\{x_n\}$ be a
Cauchy sequence of vectors in it. $V$ is said to be complete if $\{x_n\}$ 
converges to a point in $V$.
\end{defn}

\begin{defn}[Banach space]\label{s2d4}
A complete, normed vector space is called a Banach space.
\end{defn}

Not all normed-vector spaces are complete. For example,
\begin{prop}
The space $C^0[0, 1]$ with the norm defined by \eqref{s2e1} is not complete.
\end{prop}
\begin{proof}
Construct the sequence of functions $f_n$ on $[0, 1]$ as
\[
f_n(x) = \begin{cases} 0 \text{ if } x \in [0, 1/2] \\
n\left(x - \frac{1}{2}\right) \text{ if } x \in [1/2, 1/2 + 1/n] \\
1 \text{ if } x \in [1/2 + 1/n, 1].
\end{cases}
\]
Clearly, $f_n \in C^0[0, 1]$. Fix an $\epsilon > 0$. Then, for $n > m$,
\begin{eqnarray*}
\lVert f_n - f_m\rVert_1 &=& \int_0^1|f_n(x) - f_m(x)|dx \\
 &=& \int_{1/2}^{1/2+1/n}(n - m)\left(x - \frac{1}{2}\right)dx + \\
 & &  \int_{1/2 + 1/n}^{1/2 + 1/m} 1 - m\left(x - \frac{1}{2}\right)dx \\
 &=& \frac{1}{2}\left(\frac{1}{m} - \frac{1}{n}\right).
\end{eqnarray*}
If $N = \lceil 1/(2\epsilon)\rceil$ then for $n > m > N$,
\[
\lVert f_n - f_m\rVert_1 < \epsilon
\]
so that $\{f_n\}$ is a Cauchy sequence under $L^1$ norm. However, as $n
\rightarrow \infty$, $f_n$ becomes the step function
\[
f(x) = \begin{cases} 0 \text{ if } x \in [0, 1/2] \\ 1 \text { if } x \in [1/2,
1]. \end{cases}
\]
Clearly, $f$ being discontinous, does not belong to the vector space $C^0[0,1]$.
\end{proof}

\begin{defn}[Equivalence of norms]
Two norms $\lVert\cdot\rVert_a$ and $\lvert\cdot\rVert_b$ on a vector space
$V$ are said to be equivalent if there are constants $c_1, c_2 > 0$, such that
for all $v \in V$, $c_1\lVert v \rVert_b \le \lVert v \rVert_a \le c_2\lVert v
\rVert_b$.
\end{defn}

Equivalent norms define the same open sets, convergence and continuity. 

Let $V$ be a vector space of dimension $n \in \mathbb{N}$. Let $\{e_1, \ldots,
e_n\}$ be its standard basis. Then the coordinate-norm of a vector $v = 
\alpha_1e_1 + \cdots + \alpha_ne_n$ is 
\[
\lVert v \rVert_0 = \sum_{i=1}^n|\alpha_i|.
\]
We will show that any norm is equivalent to the coordinate norm through a
series of propositions.
\begin{prop}\label{s2p4}
Let $\lVert\cdot\rVert$ be an arbitrary norm on $V$ then there exists a constant
$c_1 > 0$ such that for all $v \in V$, $\lVert v \rVert \le c_1\lVert v \rVert_0
$.
\end{prop}
\begin{proof}
If $v = \sum_{i=1}^n \alpha_i e_i$, $\lVert v \rVert \le \sum_{i=1}^n\lVert 
\alpha_i e_i\rVert = \sum_{i=1}^n|\alpha_i|\lVert e_i \rVert$. Let 
\[
c_1 = \max(\{\lVert e_i \rVert\;:\; i = 1, 2, \ldots n\})
\]
then $\lVert v \rVert \le c_1\sum_{i=1}^n|\alpha_i| = c_1\lVert v \rVert_0$.
\end{proof}
\begin{rem}
We cannot assume an orthonormal basis at this stage because we do not yet have
an inner product on this space.
\end{rem}

Define a unit sphere on $V$ under the coordinate norm. It is
\[
S = \{v \in V \;:\; \lVert v \rVert_0 = 1\}.
\]

\begin{prop}\label{s2p5}
The unit sphere $S$ in the vector space $V$ under the coordinate norm is 
bounded and closed.
\end{prop}
\begin{proof}
$S$ is bounded because all its elements have a unit norm. To prove that it is
closed, we show that its complement $S^c = \{v \in V \;;\; \lVert v \rVert_0
\ne 1\}$ is open. Fix a point $u \in S^c$ and a sufficiently small $\epsilon >
0$. Define an open ball of radius $\epsilon$ centred at $u$. It is the set 
$B_\epsilon(u) = \{w \in V \;;\; d_0(u, w) < \epsilon\}$. For any $w \in 
B_\epsilon(u)$, $|\lVert w \rVert_0 - \lVert u \rVert_0| \le \lVert u - 
w\rVert_0 = d_0(u, w) < \epsilon$, by the reverse triangle inequality. 
Therefore, $-\epsilon < \lVert w \rVert_0 - \lVert u \rVert_0 < \epsilon$ or 
that $\lVert u \rVert_0 - \epsilon < \lVert w \rVert_0 < \lVert u \rVert_0 + 
\epsilon$. Since $u \in S^c$, $\lVert u \rVert_0 \ne 1$. Let $|\lVert u 
\rVert_0 - 1| = \delta$. Choose $\epsilon = \delta/2$. If $\lVert u \rVert_0 =
1 + \delta$ then we have $\lVert u \rVert_0 - \epsilon = 1 + \delta/2 < \lVert
w \rVert_0$ so that $w \in S^c$. On the other hand, if $\lVert u \rVert_0 = 1 
- \delta$ then $\lVert w \rVert_0 < 1 - \delta/2$ so that $w$ is once again
a member of $S^c$. This makes $S^c$ and $S$ closed.
\end{proof}

\begin{prop}\label{s2p6}
If $V$ is finite dimensional then the unit sphere defined under the coordinate
norm is compact.
\end{prop}
\begin{proof}
Heine-Borel theorem is applicable only to subsets of $\mathbb{R}^n$. Therefore,
for only finite dimensional vector spaces is the closed and bounded unit sphere
compact.
\end{proof}
\begin{rem}
There is a lemma attributed to F. Reisz which states that a normed vector 
space is finite dimensional if and only if the unit ball in it is compact.
\end{rem}

\begin{prop}\label{s2p7}
Any norm is a continuous function over $V$.
\end{prop}
\begin{proof}
Fix an $\epsilon > 0$ and let $|\lVert u \rVert - \lVert v \rVert| < \epsilon$.
Then $\lVert u - v\rVert < \delta = \epsilon \Rightarrow |\lVert u \rVert - 
\lVert v \rVert| < \epsilon$.
\end{proof}

\begin{prop}\label{s2p8}
For an arbitrary norm $\lVert\cdot\rVert$, there exits $c_2 > 0$ such that
$c_2\lVert u \rVert_0 \le \lVert u \rVert$.
\end{prop}
\begin{proof}
Define $f: S \rightarrow \mathbb{R}$ such that $f(v) = \lVert v \rVert$. From
propositions \ref{s2p6} and \ref{s2p7}, $f$ is a continuous function on a 
compact set. Therefore, it attains its extremal values. In particular, it
attains its minimal value, say $c_2$. That is, there exists $v^\ast \in S$
such that $f(v^\ast) = c_2$. Since $v^\ast \in S$, $\lVert v^\ast \rVert_0 \ne
0$ so that $v^\ast \ne 0$. Since $v^\ast \ne 0, f(v^\ast) = \lVert v^\ast
\rVert \ne 0$ or that $c_2 > 0$.

Let $x$ be an arbitrary vector in $V$. Then $x/\lVert x \rVert_0 \in S$ and
hence $f(x) \ge c_2$, that is,
\[
\left\lVert \frac{x}{\lVert x \rVert_0} \right\rVert \ge c_2 \Rightarrow 
\lVert x \rVert \ge c_2 \lVert x \rVert_0.
\]
\end{proof}

\begin{prop}\label{s2p9}
If $V$ is finite-dimensional then any norm is equivalent to the coordinate
norm.
\end{prop}
\begin{proof}
Follows from propositions \ref{s2p4} and \ref{s2p8}.
\end{proof}

\begin{thm}\label{s2t1}
If $V$ is finite-dimensional then all norms are equivalent.
\end{thm}
\begin{proof}
Let $\lVert\cdot\rVert_1$ and $\lVert\cdot\rVert_2$ be two norms on a finite
dimensional vector space $V$. By propostion \ref{s2p9}, there are positive
constants $c_1, c_2$ such that for any vector $v \in V$,
\begin{equation}\label{s2e2}
c_1\lVert v \rVert_0 \le \lVert v \rVert_1 \le c_2 \lVert v \rVert_0.
\end{equation}
Similarly, there are positive constants $d_1, d_2$ such that for any $v \in V$,
\begin{equation}\label{s2e3}
d_1\lVert v \rVert_0 \le \lVert v \rVert_2 \le d_2 \lVert v \rVert_0.
\end{equation}
From \eqref{s2e3}, $c_2d_1\lVert v \rVert_0 \le c_2\lVert v \rVert_2$. From
equation \eqref{s2e2}, 
\begin{equation}\label{s2e4}
d_1\lVert v \rVert_1 \le c_2\lVert v \rVert_2.
\end{equation}
Once again, from \eqref{s2e3}, we have 
\begin{equation}\label{s2e5}
c_1\lVert v \rVert_2 \le c_1d_2\lVert v \rVert_0 \le d_2\lVert v \rVert_1,
\end{equation}
where we used \eqref{s2e2} to get the last
inequality. From \eqref{s2e4} and \eqref{s2e4} we get
\[
\frac{d_1}{c_2}\lVert v \rVert_1 \le \lVert v \rVert_2 \le \frac{d_2}{c_1}
\lVert v \rVert_1.
\]
These manipulations are possible because $c_1, c_2 > 0$.
\end{proof}

The equivalence of all norms in finite dimensional spaces implies that the
topology is independent of how we measure distances. The behaviour of open sets
and convergence of seqeuences is the same for all norms. This is not so in
infinite dimensional spaces. This is why the analysis of partial differential
equations is more difficult and delicate.
\end{document}

